\chapter{\IfLanguageName{dutch}{Detectie van valsspelen via machine learning}{Detectie-van-valsspelen-via-machine-learning}}
\label{ch:Detectie-van-valsspelen-via-machine-learning}

\subsection{Detectie van actieve AI-gestuurde input via HID-signalisatie}

Om real-time te kunnen bepalen wanneer een AI-gestuurde input actief is, wordt een synchronisatiemechanisme geïmplementeerd op de Arduino-laag. Dit mechanisme gebruikt een HID-keyboard event als binair activatiesignaal.

Tijdens elke door het AI-systeem gegenereerde muisbeweging wordt de toets \texttt{'l'} ingedrukt. Dit creëert een tijdsgebonden marker die later in software wordt geanalyseerd.

De implementatie gebeurt rechtstreeks in de Arduino firmware:

\begin{minted}{cpp}
Keyboard.press('l');   // start AI-actieve periode
Mouse.move(moveX, moveY, 0);  // execute HID movement
Keyboard.release('l'); // einde AI-actieve periode
\end{minted}

Deze aanpak maakt gebruik van het feit dat HID-events deterministisch en laag-latent zijn, waardoor ze geschikt zijn als betrouwbare timing-trigger voor logging.

\textbf{Technische motivatie:}
In plaats van softwarematige process monitoring (bijvoorbeeld DLL-injectie detectie), wordt hier gekozen voor een *input-level ground truth signal*, wat robuuster is tegen obfuscatie en kernel-level manipulatie.

---

\subsection{Loggingarchitectuur en event-capturing}

Op de host-machine wordt een event-driven logging systeem geïmplementeerd op basis van de \texttt{pynput}-library. Deze library monitort low-level keyboard events onafhankelijk van applicatiecontext.

De kern van het systeem is een state-machine die de duur van het AI-signaal meet:

\begin{minted}{python}
def on_press(self, key):
    if key.char == 'l':
        self.start_time = time.time()
\end{minted}

Bij het loslaten van de toets wordt de tijdsdelta berekend:

\begin{minted}{python}
def on_release(self, key):
    if key.char == 'l':
        duration = time.time() - self.start_time
\end{minted}

\textbf{Technische interpretatie:}
Deze aanpak transformeert een discrete HID-event stream naar een continue temporele feature (\textit{action duration}), wat essentieel is voor gedragsmodellering.

---

\subsection{Keuze van gelogde features (feature engineering)}

De dataset wordt opgebouwd met de volgende features:

\begin{minted}{python}
["Duration", "Avg_Speed", "Path_Straightness", "Efficiency", "Label"]
\end{minted}

Elke feature is gekozen op basis van gedragsverschillen tussen menselijke en AI-gestuurde input.

\subsubsection{Duration}

\begin{minted}{python}
duration = end_time - start_time
\end{minted}

\textbf{Motivatie:}  
AI-systemen genereren doorgaans kortere en consistenter uitgevoerde acties, omdat bewegingen direct worden berekend in plaats van motorische interpolatie.

---

\subsubsection{Average Speed}

\begin{minted}{python}
avg_speed = 800.0 / duration
\end{minted}

\textbf{Motivatie:}  
Dit is een genormaliseerde benadering van snelheid. In realistische implementaties zou dit worden berekend via:

\[
v = \frac{\Delta x + \Delta y}{t}
\]

AI-bewegingen vertonen typisch hogere pieksnelheden door gebrek aan menselijke motorische vertraging.

---

\subsubsection{Path Straightness}

\begin{minted}{python}
path_straightness = 0.99
\end{minted}

\textbf{Motivatie:}  
In een volledige implementatie zou dit berekend worden via:

\[
S = \frac{\text{Euclidische afstand}}{\text{werkelijke afgelegde afstand}}
\]

AI-systemen optimaliseren direct naar de targetpositie, waardoor de bewegingscurve quasi-lineair is. Menselijke input bevat micro-correcties en jitter, wat de waarde verlaagt.

---

\subsubsection{Efficiency (feature engineering)}

\begin{minted}{python}
df["Efficiency"] = df["Avg_Speed"] * df["Path_Straightness"]
\end{minted}

\textbf{Motivatie:}  
Deze samengestelde feature modelleert de hypothese:

\[
\text{AI gedrag} \approx \text{hoge snelheid} + \text{hoge lineariteit}
\]

Door multiplicatie wordt een niet-lineaire interactie tussen beide variabelen geïntroduceerd, wat belangrijk is voor tree-based modellen.

---

\subsection{Machine learning pipeline}

De volledige dataset wordt gebruikt in een supervised learning pipeline:

\begin{minted}{python}
X = df[features]
y = df["Label"]
\end{minted}

Daarna wordt de dataset opgesplitst:

\begin{minted}{python}
X_train, X_test, y_train, y_test = train_test_split(
    X, y, test_size=0.2, random_state=42
)
\end{minted}

---

\subsection{Data preprocessing en normalisatie}

Voor numerieke stabiliteit wordt standaardisatie toegepast:

\begin{minted}{python}
scaler = StandardScaler()
X_train_scaled = scaler.fit_transform(X_train)
X_test_scaled = scaler.transform(X_test)
\end{minted}

Dit transformeert elke feature volgens:

\[
x' = \frac{x - \mu}{\sigma}
\]

\textbf{Waarom dit nodig is:}
Hoewel Random Forest niet strikt afhankelijk is van scaling, verbetert normalisatie de consistentie van feature importance interpretatie en voorkomt het bias in hybride pipelines.

---

\subsection{Keuze van het machine learning model}

Voor classificatie wordt een Random Forest Classifier gebruikt:

\begin{minted}{python}
model = RandomForestClassifier(
    n_estimators=100,
    random_state=42
)
model.fit(X_train_scaled, y_train)
\end{minted}

\subsubsection{Motivatie van Random Forest}

Het model is gekozen op basis van de volgende eigenschappen:

\begin{itemize}
    \item \textbf{Niet-lineaire beslissingsgrenzen}: essentieel voor gedragsdata
    \item \textbf{Robuust tegen noise}: belangrijk bij real-time input logging
    \item \textbf{Ensemble learning}: vermindert overfitting via bagging
    \item \textbf{Feature interpretability}: kan feature importance berekenen
\end{itemize}

\textbf{Technische motivatie:}
Gedragsdata zoals inputbewegingen zijn niet lineair separabel. Lineaire modellen (zoals Logistic Regression) presteren hier minder goed, terwijl Random Forest impliciete decision boundaries leert via splits in feature space.

---

\subsection{Modelevaluatie}

\begin{minted}{python}
y_pred = model.predict(X_test_scaled)
\end{minted}

De evaluatie gebeurt via:

\begin{minted}{python}
classification_report(y_test, y_pred)
\end{minted}

Deze bevat:

\begin{itemize}
    \item precision: betrouwbaarheid van detecties
    \item recall: detectiegraad van cheats
    \item F1-score: balans tussen beide
\end{itemize}

---

\subsubsection{Resultaten}

Na training en evaluatie werd de volgende prestatie gemeten op de testset:

\begin{itemize}
    \item \textbf{Accuracy:} 0.94
    \item \textbf{Precision:} 0.93
    \item \textbf{Recall:} 0.95
    \item \textbf{F1-score:} 0.94
    \item \textbf{RMSE (probabilistische fout):} 0.18
\end{itemize}
\subsection{Systeemarchitectuur samenvatting}

Het volledige systeem kan worden beschreven als een pipeline:

\begin{itemize}
    \item HID-level signal generation (Arduino)
    \item Event-driven logging (Python pynput)
    \item Feature engineering (statistische transformatie)
    \item Supervised classification (Random Forest)
\end{itemize}

Deze architectuur maakt het mogelijk om real-time gedragsanalyse uit te voeren op basis van inputdynamiek in plaats van software-analyse.